\section{Zeta-law kernels and tail partitions}
The first layer treats zeta-normalized expressions as probability, kernel, or tail-partition identities.  This viewpoint turns several inequalities into positivity statements for moments, finite shadows, and reciprocal partition functions.
\begin{proposition}[Qi-Agarwal divisor-polygamma \(f_{2}\) is completely monotone]
\label{res:qa-divisor-polygamma-f2-cm}
The second divisor-polygamma sum \(f_2(x)=[\psi'(x)]^2+\psi''(x)\) is completely monotone on \((0,\infty)\).
\end{proposition}
\begin{proof}
For every \(k\ge1\), the standard polygamma sign theorem gives
\[
(-1)^{k+1}\psi^{(k)}(x)
=
\int_0^\infty \frac{t^k e^{-xt}}{1-e^{-t}}\,dt,
\]
so \((-1)^{k+1}\psi^{(k)}\) is completely monotone on \((0,\infty)\).

If \(n\) is odd and \(km=n\), then \(k\) and \(m\) are both odd. Hence \(\psi^{(k)}\) itself is a positive completely monotone function, and
\[
\bigl[\psi^{(k)}(x)\bigr]^m
\]
is completely monotone by finite product closure of positive completely monotone functions. Summing over the finitely many divisor pairs \(km=n\) preserves complete monotonicity. Therefore \(f_n\) is completely monotone for every odd \(n\).

For \(n=2\), the divisor pairs are \((k,m)=(1,2)\) and \((2,1)\), so
\[
f_2(x)=\bigl[\psi'(x)\bigr]^2+\psi''(x).
\]
The same Qi--Agarwal source records the known sharp theorem that
\[
\bigl[\psi'(x)\bigr]^2+\lambda\psi''(x)
\]
is completely monotonic on \((0,\infty)\) if and only if \(\lambda\le1\). Taking \(\lambda=1\) gives \(f_2\in CM(0,\infty)\).

Thus the source's proposed even-parity clause, "\(f_{2\ell}\) is not completely monotonic for all \(\ell\in\mathbb N\)", is false already at \(\ell=1\). The source problem is solved negatively/correctively:

the odd half is true;
the stated even half is false;
the corrected even frontier begins at \(n=4\), not \(n=2\).
\end{proof}

\begin{lemma}[incomplete beta tail derivative has positive discrete Laplace representation for lambda...]
\label{res:incomplete-beta-tail-derivative-cm-certificate}
For \(b>0\) and \(0<\lambda\le1\), the derivative of \(I_{b,\lambda}(x)=B(b,\lambda)-B(b,\lambda;e^{-x})\) is completely monotone; for \(0<\lambda<1\), \(I_{b,\lambda}'(x)=\sum_{n=0}^\infty (1-\lambda)_n e^{-(b+n)x}/n!\).
\end{lemma}
\begin{proof}
For \(b>0\) and \(0<\lambda\le1\), define
\[
I_{b,\lambda}(x)=B(b,\lambda)-B(b,\lambda;e^{-x}).
\]
With the lower incomplete beta convention
\[
B(b,\lambda;z)=\int_0^z u^{b-1}(1-u)^{\lambda-1}\,du,
\]
one has
\[
I_{b,\lambda}(x)=\int_0^x e^{-bt}(1-e^{-t})^{\lambda-1}\,dt,
\]
and \(I_{b,\lambda}\) is a Bernstein function on \((0,\infty)\).

The substitution \(u=e^{-t}\) gives
\[
B(b,\lambda)-B(b,\lambda;e^{-x})
=\int_{e^{-x}}^1u^{b-1}(1-u)^{\lambda-1}\,du
=\int_0^x e^{-bt}(1-e^{-t})^{\lambda-1}\,dt.
\]
Therefore
\[
I_{b,\lambda}'(x)=e^{-bx}(1-e^{-x})^{\lambda-1}.
\]
For \(\lambda=1\), this is \(e^{-bx}\), a completely monotone function.

For \(0<\lambda<1\), put \(c=1-\lambda\). Then \(c\in(0,1)\), and
\[
(1-e^{-x})^{-c}
=\sum_{n=0}^\infty \frac{(c)_n}{n!}e^{-nx}
\]
with positive coefficients and locally uniform convergence on \((0,\infty)\). Thus
\[
I_{b,\lambda}'(x)
=\sum_{n=0}^\infty \frac{(c)_n}{n!}e^{-(b+n)x},
\]
the Laplace transform of the positive discrete measure
\[
\sum_{n=0}^\infty \frac{(c)_n}{n!}\delta_{b+n}.
\]
So \(I_{b,\lambda}'\) is completely monotone. Since \(I_{b,\lambda}(x)\ge0\), \(I_{b,\lambda}\in C^\infty(0,\infty)\), and \(I_{b,\lambda}'\) is completely monotone, \(I_{b,\lambda}\) is a Bernstein function.
\end{proof}

\begin{lemma}[log function signed derivative polynomial normal form and tail vanishing]
\label{res:log-function-signed-derivative-tail-vanishing}
For every \(f\in\mathcal F_{1,n}\), where \(f(x)=p(\log x)/x\), and every \(k\ge0\), the signed derivative \(L^k f\) with \(L=-d/dx\) has the form \(x^{-k-1}p_k(\log x)\) for a polynomial \(p_k\), and therefore \(L^k f(x)\to0\) as \(x\to\infty\).
\end{lemma}
\begin{proof}
We first prove that for each \(k\ge0\),
\[
L^k f(x)=x^{-k-1}p_k(\log x)
\]
for some polynomial \(p_k\).

For \(k=0\), this is the definition with \(p_0=p\).  Suppose
\[
L^k f(x)=x^{-k-1}p_k(\log x).
\]
Then
\[
\begin{aligned}
L^{k+1}f(x)
&=-\frac{d}{dx}\left(x^{-k-1}p_k(\log x)\right)\\
&=x^{-k-2}\left((k+1)p_k(\log x)-p_k'(\log x)\right).
\end{aligned}
\]
Thus \(p_{k+1}(y)=(k+1)p_k(y)-p_k'(y)\), which is again a polynomial.

Since every polynomial in \(\log x\) grows slower than every positive power of \(x\),
\[
x^{-k-1}p_k(\log x)\to0
\qquad (x\to\infty).
\]
Therefore
\[
\lim_{x\to\infty}L^k f(x)=0
\]
for every \(k\ge0\).

Assume \(f\in D_{k+1}\).  Then
\[
L^{k+1}f(x)\ge0
\]
on \((0,\infty)\).  Put
\[
g(x)=L^k f(x).
\]
By definition of \(L\),
\[
-g'(x)=L^{k+1}f(x)\ge0.
\]
Using \(g(x)\to0\) as \(x\to\infty\), integrate from \(x\) to \(R\):
\[
g(x)-g(R)=\int_x^R L^{k+1}f(t)\,dt.
\]
Letting \(R\to\infty\) gives
\[
L^k f(x)=g(x)=\int_x^\infty L^{k+1}f(t)\,dt\ge0.
\]
Hence \(f\in D_k\).  Therefore
\[
D_{k+1}\subseteq D_k
\]
for all \(k\ge0\) and all \(n\).
\end{proof}
